\documentclass[a4paper,12pt]{article}
\usepackage{graphicx}
\usepackage{hyperref}
\usepackage[italian]{babel}
\usepackage{parskip}
\setlength{\parskip}{1em}
\graphicspath{ {./images/} }
\hypersetup {
    colorlinks=true,
    linkcolor=black,
    pdfborderstyle={/S/U/W 1},
}

\begin{document}
\begin{titlepage}
    \begin{center}
        \begin{minipage}{0.3\textwidth}
            \includegraphics[width=\linewidth]{Logo_Universita_Padova300x300.png}
        \end{minipage}
        
    \end{center}

    \begin{center}
        \noindent\rule{0.7\textwidth}{0.4pt}\\[0.5cm]
        \large\textbf{Università degli Studi di Padova}\\
        \large Corso di Laurea in Informatica\\
        \small A. A. 2025/26
    \end{center}

    \vspace{1,5cm}    

    \begin{center}
        \begin{minipage}{0.5\textwidth}
            \includegraphics[width=\linewidth]{logo.png}
        \end{minipage}
    \end{center}

    \vspace{1,5cm}    
    \begin{center}
        {\Large \textbf{Relazione}\par}
        \vspace{0.6cm}
        {\Large Progetto di Tecnologie Web\par}
        \vspace{0.3cm}
    \end{center}
    \vfill

    \noindent
    \begin{tabular}{r|l}
        \textbf{Studenti} & Sofia De Blasi - matricola - email \\
         & Nenad Radulovic - matricola - email \\
         & Giacomo Speggiorin - matricola - email \\
         & Nicola Simionato - matricola - email \\
         & \\
        \textbf{Indirizzo del sito} & indirizzo\\

    \end{tabular}

    \vfill

\end{titlepage}     
\newpage
\tableofcontents
\newpage

\section{Introduzione}
\subsection{Com'è nata l'idea}
L'idea di questo progetto nasce da una domanda che ci siamo posti più volte nella nostra quotidianità: "Cosa potremmo fare questo weekend?". 
Spesso nelle città universitarie vengono organizzati eventi o attività per gli studenti. 
È però seccante doverne andare ogni volta alla ricerca. 
Bisogna sapere inanzi tutto cosa si sta cercando (il che non è una cosa scontata), e serve poi andarsi a sfogliare di volta in volta le pagine social degli organizzatori, che sono sempre diversi.\par
Abbiamo quindi pensato a quanto sarebbe comodo se fossero gli organizzatori a venire in cerca degli studenti, e a quanto si semplificherebbe la ricerca se tutti gli organizzatori avessero un'unica enorme bacheca a cui appendere i loro volantini.\par
Nel mentre buttavamo giù le idee il discorso si è ampliato, e ci siamo resi conto che la fuori ci sono un sacco di servizi di cui gli studenti hanno bisogno ma non hanno un punto di riferimento sul dove andarli a cercare.

\subsection{L'obbiettivo di StudentSpace}
StudentSpace propone una "bacheca online" che raccolga in un unico punto eventi e attività che possano interessare agli studenti al di fuori dell'ambito universitario.\par
Sulla bacheca gli studenti potranno trovare eventi sportivi, attività serali, concerti, e tanto altro per poter passare il loro tempo.
Potranno inoltre dare la loro disponibilità ad altri studenti per ripetizioni, o se necessario informarsi sugli affitti disponibili nella loro città.\par
L'obbiettivo è quello di offrire a studenti e studentesse un accesso rapido e completo a ciò che potrebbe interessarli. Inoltre, vuole aiutare gli organizzatori che si rivolgono a questo target a proporre i loro servizi a un pubblico più ampio.

\newpage
\section{Sviluppo}

\newpage
\section{Test}

\newpage
\section{Strumenti utilizzati}

\end{document}

